\uuid{YfYc}
\titre{Limite d’une somme de Riemann }

\niveau{} 				%L1, L2, L3, MPSI, MP, PCSI, PC, PSI...
\module{ Intégration } 	%Analyse, Algèbre...
\chapitre{Intégration par somme de Riemann}   			%Continuité, Groupes, Fonctions de plusieurs variables...
\sousChapitre{}			%Optimisation, Diagonalisation d'une matrice, Calcul de dérivées partielles...

\theme{}				%Fonction de répartition, Division euclidienne de polynômes, ...
\auteur{Erwan l'Haridon}
\datecreate{ 2026-02-11}
\organisation{AMSCC}			%AMSCC, Exo7, ...
\difficulte{3}			%1, 2, 3, 4 ou 5

\contenu{

\texte{ 

}

\begin{enumerate}
\item   \question{Déterminer $\displaystyle \lim_{n\to\infty}\sum_{k=1}^{2n}\frac{k}{k^2+n^2}$.}
\indication{}
\reponse{On écrit $\displaystyle \frac{k}{k^2+n^2}=\frac1n\,\frac{\frac{k}{n}}{1+\left(\frac{k}{n}\right)^2}$. Ainsi $\displaystyle \sum_{k=1}^{2n}\frac{k}{k^2+n^2}=\frac1n\sum_{k=1}^{2n} f\!\left(\frac{k}{n}\right)$ avec $f(x)=\frac{x}{1+x^2}$. Donc la limite vaut $\displaystyle \int_0^2 \frac{x}{1+x^2}\,dx=\frac12\ln(1+x^2)\Big|_0^2=\frac12\ln 5$.}

\end{enumerate}

}